\chapter{Introduction}
\label{chap:introduction}

Building upon foundational elements from the \ac{TCP}~\cite{rfc9293} and \ac{TLS}~\cite{rfc8446}, QUIC~\cite{rfc9000} is a general-purpose transport protocol developed atop the \ac{UDP}, and offers several advantages over the traditional \ac{TCP}+\ac{TLS} stack. Its key differentiators include integrated security via \ac{TLS}, reduced round trips for data transmission by consolidating connection establishment and secure channel negotiation, mitigation of \ac{TCP}'s head-of-line blocking, and enhanced extensibility through userspace implementations. These capabilities position QUIC as an ideal foundation for protocols like \ac{HTTP/3}~\cite{rfc9114}, which already accounts for 20.5\% of global HTTP traffic~\cite{cloudflare2024}. This increasing adoption also underscores the critical need for optimizing QUIC's performance. Seen in the work from Jaeger et al.~\cite{10186365}, packet I/O represents the most significant bottleneck to QUIC's maximum achievable throughput, and it highlights the unique challenges QUIC faces being a userspace protocol due to the inherent costs of context switches and there being fewer mature optimizations and offloads than \ac{TCP}.

Fortunately, \ac{XDP} offers one way to reduce packet I/O within the Linux kernel. Introduced in 2016~\cite{linux_xdp_merge_2016}, \ac{XDP} functions as a hook in Linux's network receive (RX) path, enabling packet handling at the earliest possible stage, prior to any other kernel processing. With AF\_XDP, a specific socket type, raw Layer 2 frames can be directly passed from the \ac{XDP} hook point to userspace QUIC implementations. While requiring some additional header processing in userspace, AF\_XDP allows packets to bypass the Linux network stack entirely.

Beyond maximizing throughput, many applications are highly latency-sensitive. Consequently, our research also examines the time difference for a packet to traverse from the \ac{NIC} into userspace when utilizing \ac{XDP} and MsQuic's default datapath. We use \ac{PTP} to make these measurements possible.

Finally, it is important to note that \ac{XDP} represents just one approach to mitigating the packet I/O overhead. Alternatives include \ac{DPDK}~\cite{dpdk}, a software framework providing libraries and \ac{NIC} polling-mode drivers to offload all packet processing from kernel space to userspace. Another promising alternative is in-kernel QUIC~\cite{quic_lxin}, which relocates QUIC into the kernel. In sum, each technology offers different benefits and drawbacks, which highlights the need to create a taxonomy in which we evaluate these QUIC optimizations, organizing each based on implementation complexity, performance improvement potential, unwanted side effects, flexibility, privilege level usage, and deployment.

\section{Research Questions}
\label{sec:research-questions}

In this thesis, we answer the following research questions:

\textbf{To what extent can we increase QUIC performance by using XDP in a trivial way? Is there optimization potential?} We interpret performance as both maximizing throughput and minimizing latency. MsQuic is used for this research as it sets the bar high with its optimized default datapath, which enables it to have one of the highest throughputs out of any QUIC implementation~\cite{spaeth2023}. Moreover, it is one of the few implementations to offer rudimentary support for \ac{XDP}. A thorough understanding of XDP's performance may therefore serve as a guide for other QUIC implementations that are looking to become more optimized in both throughput and latency.

\textbf{Can a taxonomy be created for QUIC optimizations, categorizing them by implementation complexity, performance improvement potential, unwanted side effects, flexibility, privilege level usage, and deployment?} As previously alluded to, there exist multiple methods to reduce packet I/O within QUIC. Each method presents its own set of tradeoffs with respect to the concepts introduced in the research question. Elaborating on each concept will give deeper insight when it is appropriate to use each technology.

\section{Outline}
\label{sec:outline}

In Chapter~\ref{chap:background}, we provide additional context on the QUIC protocol, \ac{XDP}, AF\_XDP, the \ac{QIR} used to gather throughput and latency data, MsQuic and its default datapath optimizations, and \ac{PTP}. Chapter~\ref{chap:implementation} describes the adaptations made to MsQuic and the \ac{QIR} in order to enable our measurements. We also describe what we measured and our procedure for doing so. Chapter~\ref{chap:evaluation} investigates the results of our measurements and offers explanations for them. Finally, Chapter~\ref{chap:related-work} investigates other relevant work, and Chapter~\ref{chap:conclusion} summarizes our findings and provides an outline for future work.
